\chapter{Complexity — Path Hard Machine}
%%% generated by the blueprint pipeline from TCSlib.Complexity.SpaceComplexity.Logspace.PathHardMachine
\section{Overview}

Declarations of the Lean module \texttt{TCSlib.Complexity.SpaceComplexity.Logspace.PathHardMachine}.

\section{Declarations}

\begin{definition}[Path construction steps]
\lean{Complexity.PathHard.CpS}
\label{Complexity.PathHard.CpS}
\leanok
Define a set of twelve labels for construction steps: assign $x:=p$; assign $t:=x$;
assign $t:=c$; assign $c:=e$ before setting $2n:=2^{e+1}$; assign $m:=n$; assign $c:=e$
before setting $m:=m\cdot2^e$; assign $k:=m$; assign $q:=p$; assign $u:=i$; assign
$c:=e$ before setting $u:=u/2^e$; assign $t:=u$; and assign $c:=e$ before setting
$t:=t\cdot2^e$.
\end{definition}

\begin{definition}[Shift sites]
\lean{Complexity.PathHard.ShS}
\label{Complexity.PathHard.ShS}
\leanok
The shift-site type consists of four cases: one for shifting the unary length parameter,
one for shifting the doubled adjacency-word length, one for shifting the width cap, and
one for shifting the source index of a matrix letter.
\end{definition}

\begin{definition}[Subtraction sites]
\lean{Complexity.PathHard.SbS}
\label{Complexity.PathHard.SbS}
\leanok
The subtraction-site type consists of three cases: subtracting the unary-block length
from the input position, subtracting the doubled adjacency-word length, and subtracting
the matrix-letter length to reach its target index.
\end{definition}

\begin{definition}[Two comparison labels]
\lean{Complexity.PathHard.LeS}
\label{Complexity.PathHard.LeS}
\leanok
Let $S$ be a two-element set with elements labeled $N$ and $M$.
\end{definition}

\begin{definition}[Instruction labels for the path-hardness machine]
\lean{Complexity.PathHard.Lb}
\label{Complexity.PathHard.Lb}
\leanok
\uses{Complexity.PathHard.CpS, Complexity.PathHard.LeS, Complexity.PathHard.SbS, Complexity.PathHard.ShS}
For natural numbers $a$, $b$, and $f$, define the type of instruction labels to consist
of the constructors listed below: the unindexed labels, together with labels indexed by
a Boolean value, an element of $\text{Fin }2$, $\text{Fin }3$, $\text{Fin }4$, $\text{Fin
}5$, $\text{Fin }9$, $\text{Fin }a$, $\text{Fin }b$, or $\text{Fin }f$, in exactly the
combinations specified by the constructors. Here $\text{Fin }m$ denotes the set of
natural numbers less than $m$.
\end{definition}

\begin{definition}[Copy-site destination register]
\lean{Complexity.PathHard.cpDst}
\label{Complexity.PathHard.cpDst}
\leanok
\uses{Complexity.PathHard.CpS, Complexity.PathHard.rC, Complexity.PathHard.rK, Complexity.PathHard.rM, Complexity.PathHard.rQ, Complexity.PathHard.rT, Complexity.PathHard.rU, Complexity.PathHard.rX}
The destination register of a copy site is assigned as follows: $X$ and $QP$ target
registers $2$ and $10$, respectively; $T$ and $TC$ target register $3$; $CN$, $CM$,
$CU$, and $CT$ target register $4$; $MN$ targets register $8$; $KM$ targets register
$9$; $UI$ targets register $13$; and $TU$ targets register $3$.
\end{definition}

\begin{definition}[Copy-site source register]
\lean{Complexity.PathHard.cpSrc}
\label{Complexity.PathHard.cpSrc}
\leanok
\uses{Complexity.PathHard.CpS, Complexity.PathHard.rC, Complexity.PathHard.rE, Complexity.PathHard.rI, Complexity.PathHard.rM, Complexity.PathHard.rN, Complexity.PathHard.rP, Complexity.PathHard.rU, Complexity.PathHard.rX}
The source register of a copy site is assigned as follows: $X$ and $QP$ read registers
$1$; $T$ reads register $2$; $TC$ reads register $4$; $CN$, $CM$, $CU$, and $CT$ read
register $6$; $MN$ reads register $7$; $KM$ reads register $8$; $UI$ reads register
$11$; and $TU$ reads register $13$.
\end{definition}

\begin{definition}[Next instruction at a copy site]
\lean{Complexity.PathHard.cpNext}
\label{Complexity.PathHard.cpNext}
\leanok
\uses{Complexity.PathHard.CpS, Complexity.PathHard.Lb}
For natural numbers $a$, $b$, and $f$, the next instruction associated to each copy site
is: $X$ leads to the first decrement of $x$, $T$ to the first increment of $t$, $TC$ to
the radius test, $CN$ to setting $c=e+1$, $MN$ to the copy site $CM$, $CM$ to the
initial phase of the shift site $M$, $KM$ to setting $c=1$, $QP$ to the end test with
remembered bit true, $UI$ to the copy site $CU$, $CU$ to the initial phase of halving
$u$, $TU$ to the copy site $CT$, and $CT$ to the initial phase of the shift site $T$.
\end{definition}

\begin{definition}[Shift-site register]
\lean{Complexity.PathHard.shReg}
\label{Complexity.PathHard.shReg}
\leanok
\uses{Complexity.PathHard.ShS, Complexity.PathHard.rK, Complexity.PathHard.rM, Complexity.PathHard.rN, Complexity.PathHard.rT}
The register associated with shift site $N$ is register $7$, with shift sites $M$, $K$,
and $T$ associated to registers $8$, $9$, and $3$, respectively.
\end{definition}

\begin{definition}[Shift exit labels]
\lean{Complexity.PathHard.shExit}
\label{Complexity.PathHard.shExit}
\leanok
\uses{Complexity.PathHard.Lb, Complexity.PathHard.ShS}
For natural-number parameters $a,b,f$, define the exit label for each of the four shift
sites by sending the $N$, $M$, and $K$ sites to the copy sites for $m:=n$, $k:=m$, and
$q:=p$, respectively, and the $T$ site to phase $0$ of the subtraction site that leaves
the target coordinate.
\end{definition}

\begin{definition}[Subtraction register assignment]
\lean{Complexity.PathHard.sbReg}
\label{Complexity.PathHard.sbReg}
\leanok
\uses{Complexity.PathHard.SbS, Complexity.PathHard.rM, Complexity.PathHard.rN, Complexity.PathHard.rT}
The function from the three subtraction sites to registers in $\text{Fin}(15)$ assigns
register $7$ to site $N$, register $8$ to site $M$, and register $3$ to site $T$.
\end{definition}

\begin{definition}[Subtraction exit labels]
\lean{Complexity.PathHard.sbExit}
\label{Complexity.PathHard.sbExit}
\leanok
\uses{Complexity.PathHard.Lb, Complexity.PathHard.SbS}
For natural-number parameters $a,b,f$, define the exit label at the site past the unary
block to be the first-letter separator in the walk, the exit label at the site past the
doubled adjacency word to be tail position $0$, and the exit label at the site that
leaves the target coordinate to be the dispatch test for whether the source coordinate
is zero.
\end{definition}

\begin{definition}[Register assignment for boundary labels]
\lean{Complexity.PathHard.leReg}
\label{Complexity.PathHard.leReg}
\leanok
\uses{Complexity.PathHard.LeS, Complexity.PathHard.rM, Complexity.PathHard.rN}
The two labels $N$ and $M$ are assigned registers $7$ and $8$, respectively, among the
first $15$ register indices.
\end{definition}

\begin{definition}[Positive comparison labels]
\lean{Complexity.PathHard.leYes}
\label{Complexity.PathHard.leYes}
\leanok
\uses{Complexity.PathHard.Lb, Complexity.PathHard.LeS}
For natural numbers $a,b,f$, define a map from the two comparison labels to the
corresponding machine states: the label for the unary block maps to the state that
clears the auxiliary register after that block, and the label for the doubled adjacency
word maps to the state that clears it after that word.
\end{definition}

\begin{definition}[Labels for comparison phases]
\lean{Complexity.PathHard.leNo}
\label{Complexity.PathHard.leNo}
\leanok
\uses{Complexity.PathHard.Lb, Complexity.PathHard.LeS}
For natural numbers $a$, $b$, and $f$, define a map from the two comparison labels to
the instruction labels by sending the unary-block comparison label to the instruction
that returns true, and the doubled-adjacency comparison label to the instruction that
clears the auxiliary register in the matrix phase.
\end{definition}

\begin{definition}[Pointer-advance labels]
\lean{Complexity.PathHard.advT}
\label{Complexity.PathHard.advT}
\leanok
\uses{Complexity.PathHard.Lb}
For natural-number parameters $a,b,f$, assign the four pointer-advance choices the
following labels: choice $0$ maps to the end test before the second letter with
remembered value false, choice $1$ to that end test with remembered value true, choice
$2$ to the end test before a first letter, and choice $3$ to the copy site for $x:=p$.
\end{definition}

\begin{definition}[Conditional pointer-advance labels]
\lean{Complexity.PathHard.advF}
\label{Complexity.PathHard.advF}
\leanok
\uses{Complexity.PathHard.Lb, Complexity.PathHard.advT}
For natural-number parameters $a,b,f$, a choice $k$ among $0,1,2,3$ and a natural number
$i$, define the label to be the pointer-advance label for choice $k$ when $i=0$; when
$i\ne0$, use the corresponding label from the pointer-advance assignment: the
second-letter end test with remembered value false for $k=0$, that test with remembered
value true for $k=1$, the first-letter end test for $k=2$, and the copy site for $x:=p$
for $k=3$.
\end{definition}

\begin{definition}[Radius Increment Selector]
\lean{Complexity.PathHard.sSinc}
\label{Complexity.PathHard.sSinc}
\leanok
\uses{Complexity.PathHard.Lb}
For natural numbers $a$, $b$, and $f$, define a function on natural numbers that sends
each $i<a$ to the $i$th radius-increment label, and sends each $i\ge a$ to the
radius-test label.
\end{definition}

\begin{definition}[Exponent Increment Selector]
\lean{Complexity.PathHard.sEinc}
\label{Complexity.PathHard.sEinc}
\leanok
\uses{Complexity.PathHard.Lb}
For natural numbers $a$, $b$, and $f$, define a map on natural numbers $i$ whose value
is the $i$-th exponent-increment instruction when $i<b$, and the exponent-test
instruction when $i\ge b$.
\end{definition}

\begin{definition}[Exponent offset increment selector]
\lean{Complexity.PathHard.eInc}
\label{Complexity.PathHard.eInc}
\leanok
\uses{Complexity.PathHard.Lb}
For natural numbers $a$, $b$, and $f$, and an index $i\in\mathbb{N}$, define an
instruction that selects the $i$th exponent-offset increment site when $i<f$, and
otherwise selects the copy site used before setting $c$ equal to the exponent.
\end{definition}

\begin{definition}[Pointer advancement instruction]
\lean{Complexity.PathHard.tAdv}
\label{Complexity.PathHard.tAdv}
\leanok
\uses{Complexity.PathHard.Lb}
For natural numbers $a,b,f$, a Boolean value $c$, and an index $i\in\mathbb{N}$, define
the instruction to advance the pointer, remembering $c$, if $i=0$; if $i\ne0$, define it
to test for the end of the scan, remembering $c$.
\end{definition}

\begin{definition}[Accumulator instruction selector]
\lean{Complexity.PathHard.dInc}
\label{Complexity.PathHard.dInc}
\leanok
\uses{Complexity.PathHard.Lb}
Let $a,b,f$ be natural numbers, let $c$ be a Boolean value, and let $i$ be a natural
number. Define an instruction to be the $i$th of two accumulator increments when $i<2$,
and otherwise the test of register $d$, with $c$ remembered in either case.
\end{definition}

\begin{definition}[Tail answer selection]
\lean{Complexity.PathHard.tailAns}
\label{Complexity.PathHard.tailAns}
\leanok
For natural numbers $a,b,f$, a Boolean value $m$, and $j\in\{0,1,2,3,4\}$, define the
answer to be true if $j=1$ or $j\ge3$, and to be $m$ otherwise.
\end{definition}

\begin{definition}[Copy-site instruction]
\lean{Complexity.PathHard.cpI}
\label{Complexity.PathHard.cpI}
\leanok
\uses{Complexity.LogProg.Ins, Complexity.PathHard.CpS, Complexity.PathHard.Lb, Complexity.PathHard.cpDst, Complexity.PathHard.cpNext, Complexity.PathHard.cpSrc}
For natural numbers $a,b,f$, a copy site $k$, and a Boolean value $b$, define an
instruction that increments the destination register of $k$ and returns to $k$ when $b$
is true. When $b$ is false, it tests whether the destination and source registers of $k$
are equal, continuing to the next site if they are and proceeding to the
copy-and-increment site otherwise.
\end{definition}

\begin{definition}[Shift-site instruction]
\lean{Complexity.PathHard.shI}
\label{Complexity.PathHard.shI}
\leanok
\uses{Complexity.LogProg.Ins, Complexity.PathHard.Lb, Complexity.PathHard.rC, Complexity.PathHard.rZ}
For natural numbers $a,b,f$, a register $r$, an exit label $e$, a map $m$ from the nine
phases to labels, and a phase $j\in\{0,1,\ldots,8\}$, define the instruction as follows:
phase $0$ tests whether the designated counter register is zero, going to $e$ if so and
to phase $1$ otherwise; phase $1$ decrements that counter and goes to phase $2$. Phase
$2$ tests whether $r$ is zero, going to phase $6$ if so and to phase $3$ otherwise;
phase $3$ decrements $r$ and goes to phase $4$; phases $4$ and $5$ increment the
designated register and go to phases $5$ and $2$, respectively. Phase $6$ tests whether
the designated register is zero, going to phase $0$ if so and to phase $7$ otherwise;
phase $7$ decrements it and goes to phase $8$; phase $8$ increments $r$ and goes to
phase $6$.
\end{definition}

\begin{definition}[Subtraction instruction phases]
\lean{Complexity.PathHard.sbI}
\label{Complexity.PathHard.sbI}
\leanok
\uses{Complexity.LogProg.Ins, Complexity.PathHard.Lb, Complexity.PathHard.rI}
For natural numbers $a,b,f$, a register $q$ among fifteen registers, a label $e$, a
function $m$ assigning a label to each of three phases, and a phase $j$, define an
instruction as follows: at phase $0$, test whether $q$ is zero, continuing to $e$ if so
and to phase $1$ otherwise; at phase $1$, decrement $q$ and continue to phase $2$; at
phase $2$, decrement register $11$ and return to phase $0$.
\end{definition}

\begin{definition}[Comparison instruction phases]
\lean{Complexity.PathHard.leI}
\label{Complexity.PathHard.leI}
\leanok
\uses{Complexity.LogProg.Ins, Complexity.PathHard.Lb, Complexity.PathHard.rD, Complexity.PathHard.rI}
For natural numbers $a,b,f$, a register $r$ among fifteen registers, labels $y,n$, a
function $m$ assigning a label to each of three phases, and a phase $j$, define an
instruction as follows: at phase $0$, compare register $12$ with $r$, continuing to $y$
if they are equal and to phase $1$ otherwise; at phase $1$, compare register $12$ with
register $11$, continuing to $n$ if they are equal and to phase $2$ otherwise; at phase
$2$, increment register $12$ and return to phase $0$.
\end{definition}

\begin{definition}[Halving instruction phases]
\lean{Complexity.PathHard.hvI}
\label{Complexity.PathHard.hvI}
\leanok
\uses{Complexity.LogProg.Ins, Complexity.PathHard.Lb, Complexity.PathHard.rC, Complexity.PathHard.rU}
For natural numbers $a,b,f$ and a phase $j$ among three phases, define an instruction as
follows: at phase $0$, test whether register $4$ is zero, continuing to the copy site
that transfers register $13$ to register $3$ if so and to phase $1$ otherwise; at phase
$1$, decrement register $4$ and continue to phase $2$; at phase $2$, halve register $13$
and return to phase $0$.
\end{definition}

\begin{definition}[Path-hard machine transition function]
\lean{Complexity.PathHard.A}
\label{Complexity.PathHard.A}
\leanok
\uses{Complexity.LogProg.ARM, Complexity.PathHard.Lb, Complexity.PathHard.advF, Complexity.PathHard.cpI, Complexity.PathHard.dInc, Complexity.PathHard.eInc, Complexity.PathHard.hvI, Complexity.PathHard.leI, Complexity.PathHard.leNo, Complexity.PathHard.leReg, Complexity.PathHard.leYes, Complexity.PathHard.rC, Complexity.PathHard.rD, Complexity.PathHard.rE, Complexity.PathHard.rI, Complexity.PathHard.rK, Complexity.PathHard.rN, Complexity.PathHard.rP, Complexity.PathHard.rQ, Complexity.PathHard.rS, Complexity.PathHard.rT, Complexity.PathHard.rU, Complexity.PathHard.rX, Complexity.PathHard.rY, Complexity.PathHard.sEinc, Complexity.PathHard.sSinc, Complexity.PathHard.sbExit, Complexity.PathHard.sbI, Complexity.PathHard.sbReg, Complexity.PathHard.shExit, Complexity.PathHard.shI, Complexity.PathHard.shReg, Complexity.PathHard.tAdv, Complexity.PathHard.tailAns}
For natural numbers $a,b,f$ and a Boolean mode, define a transition function assigning a
register-machine instruction to every label of the machine. Its instructions implement
input-length measurement, scanning and decoding the input, parameter and width
calculations, and the subsequent dispatch and decider calls; the specified label cases
determine the exact instruction and continuation labels, including the copy, shift,
subtraction, comparison, and halving routines.
\end{definition}

\begin{lemma}[Program configurations satisfy the precondition]
\lean{Complexity.PathHard.preS_A}
\label{Complexity.PathHard.preS_A}
\leanok
\uses{Complexity.LogProg.AConf, Complexity.LogProg.PreS, Complexity.PathHard.A, Complexity.PathHard.Lb, Complexity.PathHard.cpDst, Complexity.PathHard.cpI, Complexity.PathHard.cpSrc, Complexity.PathHard.hvI, Complexity.PathHard.leI, Complexity.PathHard.leReg, Complexity.PathHard.sbI, Complexity.PathHard.shI}
\difficulty{5}
For any natural numbers $a,b,f$, Boolean mode $m$, oracle $o$ on five indices and finite
Boolean strings, finite Boolean string $y$, and abstract configuration $c$ with fifteen
natural-number registers, $c$ satisfies the program precondition. Thus, when its
instruction compares two registers, their indices are distinct; when it compares a
register with the input payload, $y$ is a valid plain encoding; when it compares a
register with either component of a paired payload, $y$ is a valid paired encoding; and
when it calls an oracle, its argument registers are pairwise distinct.
\end{lemma}

\begin{lemma}[Register bounds imply the guard]
\lean{Complexity.PathHard.gb_of}
\label{Complexity.PathHard.gb_of}
\leanok
\uses{Complexity.LogProg.GuardLe, Complexity.PathHard.A, Complexity.PathHard.Lb, Complexity.PathHard.St, Complexity.PathHard.St.Bd, Complexity.PathHard.St.Bd.toV, Complexity.PathHard.St.toV, Complexity.PathHard.preS_A}
\difficulty{1}
Let $a,b,f,B$ be natural numbers, let $m$ be a Boolean mode, let $o$ be an oracle on
five deciders and bit strings, and let $y$ be a bit string. If a fifteen-register state
$st$ has every register value at most $B$, then for every program label $l$, the
abstract configuration with label $l$, register values given by $st$, and no Boolean
result satisfies the guard: its registers are all at most $B$ and its
instruction-specific precondition holds.
\end{lemma}

\begin{lemma}[Removing an Initial Program Step]
\lean{Complexity.PathHard.step_of}
\label{Complexity.PathHard.step_of}
\leanok
\uses{Complexity.LogProg.AConf, Complexity.LogProg.AReach, Complexity.LogProg.AReach.step, Complexity.LogProg.GuardLe, Complexity.LogProg.astep, Complexity.PathHard.A, Complexity.PathHard.Lb, Complexity.PathHard.St, Complexity.PathHard.St.Bd, Complexity.PathHard.St.toV, Complexity.PathHard.gb_of}
\difficulty{1}
Let $a,b,f,B$ be natural numbers, let $m$ be a Boolean mode, let
$o:\text{Fin}(5)\to\text{List}(\text{Bool})\to\text{Bool}$ be an oracle, let $y$ be a Boolean
string, and let $st$ be a fifteen-register state with every register at most $B$. For
any instruction $l$ and abstract configuration $c$, if there is a finite run from the
configuration obtained by taking one step of the program at $l$ with registers $st$ and
no result, to $c$, with the guard satisfied at every configuration before the endpoint,
then there is such a guarded finite run from the configuration at $l$ with registers
$st$ and no result to $c$.
\end{lemma}

\begin{lemma}[Removing an initial program step]
\lean{Complexity.PathHard.stepH_of}
\label{Complexity.PathHard.stepH_of}
\leanok
\uses{Complexity.LogProg.AHalt, Complexity.LogProg.AHalt.step, Complexity.LogProg.GuardLe, Complexity.LogProg.astep, Complexity.PathHard.A, Complexity.PathHard.Lb, Complexity.PathHard.St, Complexity.PathHard.St.Bd, Complexity.PathHard.St.toV, Complexity.PathHard.gb_of}
\difficulty{1}
Let $a,b,f,B$ be natural numbers, let $mode$ be Boolean, let $o$ be an oracle assigning
a Boolean value to each of five indices and each Boolean word, let $y$ be a Boolean
word, and let $st$ be a state of fifteen natural-number registers, each at most $B$. For
any program instruction $l$ and Boolean answer $r$, if there is a finite guarded run
that halts with answer $r$ from the configuration obtained by taking one program step
from instruction $l$ and register state $st$, then there is a finite guarded run that
halts with answer $r$ starting directly from instruction $l$ and register state $st$.
\end{lemma}

\begin{lemma}[Halting from an answer state]
\lean{Complexity.PathHard.ans_halt}
\label{Complexity.PathHard.ans_halt}
\leanok
\uses{Complexity.LogProg.AHalt, Complexity.LogProg.AHalt.ret, Complexity.LogProg.GuardLe, Complexity.PathHard.A, Complexity.PathHard.St, Complexity.PathHard.St.Bd, Complexity.PathHard.St.toV, Complexity.PathHard.gb_of}
\difficulty{1}
Let $a,b,f,B\in\mathbb{N}$, let $m\in\{\mathrm{false},\mathrm{true}\}$, let $o$ be an
oracle on five Boolean words, and let $y$ be a Boolean word. For any register record
$st$ whose fifteen entries are at most $B$ and either Boolean $r$, the machine
$A(a,b,f,m)$ on input $y$, started at its answer-$r$ state with register values $st$ and
no result recorded, halts with result $r$; at every earlier time its configuration
satisfies the guard requiring its prescribed precondition and all register values to be
at most $B$.
\end{lemma}

\begin{lemma}[Continuation from equal register records]
\lean{Complexity.PathHard.cont}
\label{Complexity.PathHard.cont}
\leanok
\uses{Complexity.LogProg.AConf, Complexity.LogProg.AReach, Complexity.LogProg.GuardLe, Complexity.PathHard.A, Complexity.PathHard.Lb, Complexity.PathHard.St, Complexity.PathHard.St.toV}
\difficulty{1}
Let $a,b,f,B\in\mathbb{N}$, let $m$ be a Boolean mode, let $o$ be an oracle on five
indices and Boolean words, and let $y$ be a Boolean word. Let $st,st'$ be register
records, $l$ a machine label, and $c$ an abstract machine configuration. If $st=st'$ and
the machine $A(a,b,f,m)$ reaches $c$ from label $l$ with register values $st'$, no
result recorded, and the guard satisfied at every earlier time, then it also reaches $c$
from label $l$ with register values $st$, no result recorded, and the same guard
satisfied at every earlier time.
\end{lemma}

\begin{lemma}[Halting Invariance Under Equal Registers]
\lean{Complexity.PathHard.contH}
\label{Complexity.PathHard.contH}
\leanok
\uses{Complexity.LogProg.AHalt, Complexity.LogProg.GuardLe, Complexity.PathHard.A, Complexity.PathHard.Lb, Complexity.PathHard.St, Complexity.PathHard.St.toV}
\difficulty{1}
Let $a,b,f,B$ be natural numbers, let $m$ be a Boolean mode, let $o$ be an oracle taking
an index from a five-element set and a Boolean word, and let $y$ be a Boolean word. For
two register records $st,st'$, a machine label $l$, and a Boolean answer $r$, if
$st=st'$ and the machine halts with answer $r$ from label $l$ and register values given
by $st'$, starting with no answer recorded and satisfying the guard at every earlier
step, then it also halts with answer $r$ from label $l$ and register values given by
$st$, under the same guard.
\end{lemma}

\begin{lemma}[Distinct Source and Destination Registers]
\lean{Complexity.PathHard.cp_ne}
\label{Complexity.PathHard.cp_ne}
\leanok
\uses{Complexity.PathHard.CpS, Complexity.PathHard.cpDst, Complexity.PathHard.cpSrc}
\difficulty{4}
For every copy operation $k$ in the finite set of twelve copy operations, the
destination register $\text{cpDst}(k)$ differs from the source register $\text{cpSrc}(k)$.
\end{lemma}

\begin{theorem}[Execution of a Copy Site]
\lean{Complexity.PathHard.cp_run}
\label{Complexity.PathHard.cp_run}
\leanok
\uses{Complexity.LogProg.AReach, Complexity.LogProg.GuardLe, Complexity.LogProg.registerCopy_loop, Complexity.PathHard.A, Complexity.PathHard.CpS, Complexity.PathHard.St, Complexity.PathHard.St.Bd, Complexity.PathHard.St.Bd.toV, Complexity.PathHard.St.toV, Complexity.PathHard.cpDst, Complexity.PathHard.cpNext, Complexity.PathHard.cpSrc, Complexity.PathHard.cp_ne}
\difficulty{1}
Let $a,b,f,B$ be natural numbers, let $m$ be a Boolean mode, let $o$ be an oracle taking
an index from a five-element set and a Boolean word, and let $y$ be a Boolean word. For
a copy-site label $k$ and a register record $st$, suppose every register value is at
most $B$ and the value in the destination register specified by $k$ is at most the value
in its source register. Then the machine reaches the next label specified by $k$ from
the copy-site label, starting with register values $st$ and no answer recorded; the
final register vector is obtained by replacing the destination value with the source
value and leaving every other value unchanged, and the guard holds at every step before
the endpoint.
\end{theorem}

\begin{lemma}[Distinct shift registers]
\lean{Complexity.PathHard.sh_ne}
\label{Complexity.PathHard.sh_ne}
\leanok
\uses{Complexity.PathHard.ShS, Complexity.PathHard.rC, Complexity.PathHard.rZ, Complexity.PathHard.shReg}
\difficulty{4}
For each of the four shift operations, the counter register, the register being shifted,
and the scratch register are pairwise distinct.
\end{lemma}

\begin{theorem}[Bounded shift run]
\lean{Complexity.PathHard.sh_run}
\label{Complexity.PathHard.sh_run}
\leanok
\uses{Complexity.LogProg.AReach, Complexity.LogProg.GuardLe, Complexity.LogProg.shl_run, Complexity.PathHard.A, Complexity.PathHard.ShS, Complexity.PathHard.St, Complexity.PathHard.St.Bd, Complexity.PathHard.St.Bd.toV, Complexity.PathHard.St.toV, Complexity.PathHard.rC, Complexity.PathHard.shExit, Complexity.PathHard.shReg, Complexity.PathHard.sh_ne}
\difficulty{1}
Let $a,b,f,B$ be natural numbers, let $mode$ be Boolean, let $o$ be an oracle assigning
a Boolean value to each of five indices and each Boolean word, let $y$ be a Boolean
word, and let $st$ be a state of fifteen natural-number registers. Suppose every
register of $st$ is at most $B$, its shift scratch register is zero, and the register
selected by the shift operation $k$, multiplied by $2^{c}$ where $c$ is the loop-counter
register of $st$, is at most $B$. Then the register program has a finite run from phase
zero of shift operation $k$ with register values $st$ to its exit phase, with the
loop-counter register cleared and the selected register replaced by its original value
times $2^{c}$; all other registers retain their original values. The guard holds at
every configuration before the endpoint.
\end{theorem}

\begin{lemma}[Subtraction registers avoid the index register]
\lean{Complexity.PathHard.sb_ne}
\label{Complexity.PathHard.sb_ne}
\leanok
\uses{Complexity.PathHard.SbS, Complexity.PathHard.rI, Complexity.PathHard.sbReg}
\difficulty{4}
For each of the three subtraction sites, the register assigned to that site is different
from the register used for the index.
\end{lemma}

\begin{theorem}[Execution of a subtraction routine]
\lean{Complexity.PathHard.sb_run}
\label{Complexity.PathHard.sb_run}
\leanok
\uses{Complexity.LogProg.AReach, Complexity.LogProg.GuardLe, Complexity.LogProg.sub_run, Complexity.PathHard.A, Complexity.PathHard.SbS, Complexity.PathHard.St, Complexity.PathHard.St.Bd, Complexity.PathHard.St.Bd.toV, Complexity.PathHard.St.toV, Complexity.PathHard.rI, Complexity.PathHard.sbExit, Complexity.PathHard.sbReg, Complexity.PathHard.sb_ne}
\difficulty{1}
Let $a,b,f,B\in\mathbb{N}$, let $m\in\{\mathrm{false},\mathrm{true}\}$, let $o$ assign a
Boolean to each of five indices and each finite Boolean string, and let $y$ be a finite
Boolean string. For any subtraction site $k$ and register configuration $st$ whose
fifteen registers are at most $B$, the register program starting at the phase-zero
instruction for $k$, with register values given by $st$ and no answer recorded, reaches
its exit instruction. At the endpoint, the register associated with $k$ is zero and the
index register has value $\max(st.i-st_{r},0)$, where $st_r$ is the initial value of the
register associated with $k$; every configuration before the endpoint satisfies the
program's preconditions and has all register values at most $B$.
\end{theorem}

\begin{lemma}[Register separation in comparisons]
\lean{Complexity.PathHard.le_ne}
\label{Complexity.PathHard.le_ne}
\leanok
\uses{Complexity.PathHard.LeS, Complexity.PathHard.leReg, Complexity.PathHard.rD}
\difficulty{4}
For either of the two comparison phases, the scratch register is distinct from the
register representing the unary-block length or from the register representing the
adjacency-word length, respectively.
\end{lemma}

\begin{theorem}[Bounded comparison routine]
\lean{Complexity.PathHard.le_run}
\label{Complexity.PathHard.le_run}
\leanok
\uses{Complexity.Codec.le_loop, Complexity.LogProg.AReach, Complexity.LogProg.GuardLe, Complexity.PathHard.A, Complexity.PathHard.LeS, Complexity.PathHard.St, Complexity.PathHard.St.Bd, Complexity.PathHard.St.Bd.toV, Complexity.PathHard.St.toV, Complexity.PathHard.leNo, Complexity.PathHard.leReg, Complexity.PathHard.leYes, Complexity.PathHard.le_ne, Complexity.PathHard.rD, Complexity.PathHard.rI}
\difficulty{1}
Let $a,b,f,B$ be natural numbers, let $m$ be a Boolean mode, let $o$ be an oracle on
five deciders and bit strings, let $y$ be a bit string, and let $k$ be one of the two
comparison labels. Suppose a fifteen-register state $st$ has all register values at most
$B$ and its decoding-and-comparison scratch register is $0$. Starting with label the
comparison routine for $k$ in phase $0$, register values $st$, and no result, the
program reaches a configuration whose label is the success continuation if the
comparison register selected by $k$ is at most the index register $i$, and otherwise is
the failure continuation; the destination scratch register equals the minimum of that
comparison-register value and $i$, with all other registers unchanged and no result.
Every configuration strictly before the reached configuration satisfies the guard
requiring all registers to be at most $B$ and the instruction-specific preconditions to
hold.
\end{theorem}

\begin{theorem}[Halving the source register]
\lean{Complexity.PathHard.hv_run}
\label{Complexity.PathHard.hv_run}
\leanok
\uses{Complexity.LogProg.AReach, Complexity.LogProg.GuardLe, Complexity.LogProg.halve_run, Complexity.PathHard.A, Complexity.PathHard.St, Complexity.PathHard.St.Bd, Complexity.PathHard.St.Bd.toV, Complexity.PathHard.St.toV, Complexity.PathHard.rC, Complexity.PathHard.rU}
\difficulty{3}
Let $a,b,f,B$ be natural numbers, let $m$ be a Boolean mode, let $o$ be an oracle on
five deciders and bit strings, and let $y$ be a bit string. For any fifteen-register
state $st$ whose register values are all at most $B$, the program reaches a
configuration from the initial one with label the halving routine in phase $0$, register
values $st$, and no result, whose label is the copy routine transferring the source
register to the scratch register, whose loop counter is $0$, whose source register is
the original source value divided by $2$ raised to the original loop-counter value, and
whose other registers agree with $st$. Every configuration strictly before the reached
configuration satisfies the guard requiring all registers to be at most $B$ and the
instruction-specific preconditions to hold.
\end{theorem}
